% ============================================================
% 理论力学 · 交互讲义 —— M6《碰撞与变质量》note.tex
% ============================================================
\documentclass[UTF8,a4paper,11pt]{ctexart}
\usepackage{amsmath,amssymb}
\usepackage{geometry}
\usepackage{booktabs}
\usepackage{hyperref}
\geometry{margin=2.5cm}
\hypersetup{colorlinks=true,linkcolor=blue,urlcolor=blue}

\title{\textbf{碰撞与变质量}}
\author{理论力学 · 交互讲义}
\date{\today}

\begin{document}
\maketitle

\begin{abstract}
本讲义处理碰撞与变质量系统：光滑球碰撞的散射角与恢复系数、
雨滴下落的变质量运动、弹性碰撞链的反跳高度。
全部解析结果由配套脚本 \texttt{qa/numeric\_check.py} 逐项数值验证。
本模块覆盖作业-5 Q1--Q3。
\end{abstract}

\tableofcontents
\newpage

% ============================================================
\section{作业映射}
\label{sec:map}

\begin{center}
\small
\begin{tabular}{cp{6.5cm}p{3.5cm}c}
\toprule
来源 & 题面简述 & 核心知识点 & 对应小节 \\
\midrule
作业-5 Q1 & 光滑球碰撞，求偏转角 $\beta$ 最大值及对应 $\alpha$ &
  碰撞、恢复系数、动量守恒 & \ref{sub:collision} \\
作业-5 Q2 & 雨滴下落，质量增加率 $\propto$ 表面积，求 $v(t)$ &
  变质量系统、牛顿第二定律 & \ref{sub:raindrop} \\
作业-5 Q3 & 两球弹性碰撞链，求 B2 反弹高度 &
  弹性碰撞、能量与动量守恒 & \ref{sub:chain} \\
\bottomrule
\end{tabular}
\end{center}

% ============================================================
\section{模型与假设}
\label{sec:model}

\begin{enumerate}
  \item \textbf{碰撞假设}：碰撞时间极短，碰撞期间重力等外力忽略。
  \item \textbf{光滑球}：切向无摩擦，仅法向有恢复系数 $e$。
  \item \textbf{变质量}：雨滴质量随时间增加，增加的质量初速度为 0
        （静止雨水），系统为开放系统。
  \item \textbf{弹性碰撞}：Q3 中所有碰撞为完全弹性（$e=1$）。
  \item \textbf{单位制}：$g$ 为重力加速度。
\end{enumerate}

% ============================================================
\section{推导}
\label{sec:derive}

\subsection{光滑球碰撞（作业-5 Q1）}
\label{sub:collision}

质量相等的光滑球，球 1 以速度 $v$ 入射，球 2 静止。
入射方向与两球中心连线夹角为 $\alpha$。恢复系数 $e$。

分解为法向（沿中心连线）与切向：
\begin{equation}
  v_{1n}=v\cos\alpha,\quad v_{1t}=v\sin\alpha,\quad v_{2n}=v_{2t}=0.
\end{equation}
碰撞后（等质量，恢复系数 $e$）：
\begin{equation}
  v_{1n}'=\frac{1-e}{2}v\cos\alpha,\quad
  v_{2n}'=\frac{1+e}{2}v\cos\alpha,\quad
  v_{1t}'=v\sin\alpha,\quad v_{2t}'=0.
\end{equation}
球 1 偏转角 $\beta$（与原方向夹角）：
\begin{equation}
  \cos\beta=\frac{v_{1n}v_{1n}'+v_{1t}v_{1t}'}{v\cdot v_1'}
  =\frac{\frac{1-e}{2}\cos^2\alpha+\sin^2\alpha}
         {\sqrt{\left(\frac{1-e}{2}\cos\alpha\right)^2+\sin^2\alpha}}.
\end{equation}
$\beta$ 最大时满足 $\mathrm d\beta/\mathrm d\alpha=0$，解得
\begin{equation}
  \tan^2\alpha=\frac{1-e}{2}.
  \label{eq:alpha-max}
\end{equation}
代入得 $\beta_{\max}$。

\subsection{雨滴变质量运动（作业-5 Q2）}
\label{sub:raindrop}

雨滴半径 $r$，质量 $m=\rho\frac43\pi r^3$。质量增加率与表面积成正比：
\begin{equation}
  \frac{\mathrm dm}{\mathrm dt}=k\cdot 4\pi r^2.
\end{equation}
另一方面 $\mathrm dm/\mathrm dt=\rho\cdot 4\pi r^2\,\dot r$，故 $\dot r=k/\rho=\lambda$（常数）。
所以 $r(t)=r_0+\lambda t$。

变质量牛顿方程（增加质量初速度为 0）：
\begin{equation}
  \frac{\mathrm d(mv)}{\mathrm dt}=mg
  \quad\Longrightarrow\quad
  m\dot v+v\dot m=mg
  \quad\Longrightarrow\quad
  \dot v=g-\frac{v\dot m}{m}.
\end{equation}
$\dot m/m=3\lambda/(r_0+\lambda t)$，故
\begin{equation}
  \dot v+\frac{3\lambda}{r_0+\lambda t}v=g.
\end{equation}
积分因子 $\mu=(r_0+\lambda t)^3$：
\begin{equation}
  \frac{\mathrm d}{\mathrm dt}\left[v(r_0+\lambda t)^3\right]=g(r_0+\lambda t)^3.
\end{equation}
积分（$v(0)=0$）：
\begin{equation}
  \boxed{\;v(t)=\frac{g}{4\lambda}\left[(r_0+\lambda t)-\frac{r_0^4}{(r_0+\lambda t)^3}\right]\;}
  \label{eq:raindrop-v}
\end{equation}
当 $t\to\infty$，$v\to gt/4$（线性增长）。

\subsection{弹性碰撞链（作业-5 Q3）}
\label{sub:chain}

球 B1 质量 $m_1$，B2 质量 $m_2$（$m_2\ll m_1$）。B1 从高度 $h$ 落下，
与地面弹性碰撞后上升速度 $v_1=\sqrt{2gh}$，再与 B2（静止）弹性碰撞。

一维弹性碰撞公式：
\begin{equation}
  v_2'=\frac{2m_1}{m_1+m_2}v_1\approx 2v_1 \quad (m_2\ll m_1).
\end{equation}
B2 反弹高度
\begin{equation}
  H=\frac{v_2'^2}{2g}\approx\frac{4v_1^2}{2g}=4h.
\end{equation}
\begin{equation}
  \boxed{\;H\approx 4h\;}
\end{equation}

% ============================================================
\section{解析结果}
\label{sec:result}

\begin{enumerate}
  \item 球碰撞偏转角：$\cos\beta$ 公式（\eqref{sub:collision}），
        $\beta_{\max}$ 时 $\tan^2\alpha=(1-e)/2$（\eqref{eq:alpha-max}）。
  \item 雨滴速度：$v(t)=\frac{g}{4\lambda}[(r_0+\lambda t)-r_0^4/(r_0+\lambda t)^3]$
        （\eqref{eq:raindrop-v}），$t\to\infty$ 时 $v\sim gt/4$。
  \item 弹性碰撞链：$H\approx 4h$（$m_2\ll m_1$）。
\end{enumerate}

% ============================================================
\section{数值验证}
\label{sec:numeric}

全部数值实验由配套脚本 \texttt{qa/numeric\_check.py} 完成（仅依赖 numpy）。

\begin{table}[htbp]
\centering
\small
\begin{tabular}{p{5cm}cc}
\toprule
物理量 & 解析值 & 数值值 \\
\midrule
$\beta_{\max}$（$e=0.5$） & $36.9^\circ$ & $36.8699^\circ$ \\
$v(t)$（雨滴，$t=10$） & 公式值 & 数值积分一致 \\
$H/h$（碰撞链） & $4.0$ & $4.0$ \\
\bottomrule
\end{tabular}
\caption{数值验证汇总。全部数值取自 \texttt{qa/numeric\_check.py} 的实际输出。}
\label{tab:numeric}
\end{table}

% ============================================================
\subsection{碰撞的几何解释}
\label{sub:collision-geometry}

等质量光滑球碰撞：入射球偏转角 $\beta$ 依赖于碰撞参数 $\alpha$ 与恢复系数 $e$。
$e=1$（完全弹性）时 $\beta=90^\circ-\alpha$。$e=0$ 时“完全非弹性”只表示法向相对速度为零，
切向速度仍为 $v\sin\alpha$，因此一般 $\beta\ne0$：
\begin{equation}
  \boldsymbol v_1'=\left(\frac{v\cos\alpha}{2},\,v\sin\alpha\right),
  \qquad
  \beta=\arctan\!\left(2\tan\alpha\right)-\alpha.
\end{equation}
$\alpha=45^\circ$ 时 $\beta=18.435^\circ$；$e=0$ 的最大偏转为
$\beta_{\max}=19.471^\circ$（$\tan^2\alpha=1/2$），只有 $\alpha=0$ 才是 $\beta=0$。
若两球真的粘在一起，必须额外加入切向摩擦或粘接，这与“光滑球”假设矛盾。
中间 $0<e<1$ 时 $\beta$ 有最大值，对应 $\tan^2\alpha=(1-e)/2$。

\subsection{变质量系统的通用公式}
\label{sub:variable-mass-general}

一般变质量系统：$\mathrm d(m\boldsymbol v)/\mathrm dt=\boldsymbol F+\boldsymbol u\,\mathrm dm/\mathrm dt$，
其中 $\boldsymbol u$ 为增加质量的绝对速度。雨滴问题中 $\boldsymbol u=0$（静止雨水）。

\subsection{碰撞链的能量传递}
\label{sub:energy-transfer}

弹性碰撞链中，能量从重球传递到轻球。$m_2\ll m_1$ 时，
$v_2'\approx 2v_1$，能量传递效率 $\sim 4m_2/m_1$（小）。多次碰撞可逐步加速轻球。

\section{碰撞的分类}
\label{sec:collision-types}

\begin{enumerate}
  \item \textbf{完全弹性}（$e=1$）：动能守恒，相对速度反向。
  \item \textbf{完全非弹性}（$e=0$）：法向相对速度为零，动能损失最大；光滑球切向速度仍保留，一般不等于“粘在一起”。
  \item \textbf{部分弹性}（$0<e<1$）：一般情况，动能部分损失。
\end{enumerate}

\paragraph{碰撞的冲量。}
碰撞冲量 $\boldsymbol J=\int \boldsymbol F\,\mathrm dt=\Delta\boldsymbol p$。
对光滑球，冲量沿法向，切向无冲量。

\section{碰撞与变质量的应用}
\label{sec:applications}

\paragraph{台球游戏。}
台球碰撞是完全弹性碰撞（$e\approx1$）的日常例子。
碰撞角、偏转角的计算可预测台球轨迹。

\paragraph{火箭推进。}
火箭是变质量系统（$dm/dt<0$），推力 $F=-\dot m u$。
Tsiolkovsky 公式：$\Delta v=u\ln(m_0/m_f)$。

\paragraph{雨滴增长。}
雨滴下落时质量增长，速度增长慢于自由落体（$v\sim gt/4$ vs $gt$）。

\subsection{雨滴 ODE 的详细求解}
\label{sub:raindrop-detail}

方程 $\dot v+3\lambda v/(r_0+\lambda t)=g$ 的积分因子
$\mu(t)=\exp(\int 3\lambda/(r_0+\lambda t)\,\mathrm dt)=(r_0+\lambda t)^3$。
乘以 $\mu$：$\mathrm d/\mathrm dt[v(r_0+\lambda t)^3]=g(r_0+\lambda t)^3$。
积分：$v(r_0+\lambda t)^3 = g(r_0+\lambda t)^4/(4\lambda)+C$。
初值 $v(0)=0$：$C=-g r_0^4/(4\lambda)$。
最终：$v=g/(4\lambda)[(r_0+\lambda t)-r_0^4/(r_0+\lambda t)^3]$。

\subsection{碰撞链的详细分析}
\label{sub:chain-detail}

B1 从高度 $h$ 落下：$v_1=\sqrt{2gh}$（向下）。
与地面弹性碰撞：$v_1'=-v_1$（向上）。
与 B2 弹性碰撞（$m_2\ll m_1$）：$v_2'=2v_1$（向上）。
B2 上升高度 $H=v_2'^2/(2g)=2v_1^2/g=4h$。

\section{附录：碰撞公式汇总}
\label{sec:appendix}

\begin{enumerate}
  \item 一维弹性碰撞：$v_1'=(m_1-m_2)v_1/(m_1+m_2)$，$v_2'=2m_1 v_1/(m_1+m_2)$。
  \item 二维光滑球碰撞：法向恢复系数 $e$，切向不变。
  \item 变质量方程：$\mathrm d(mv)/\mathrm dt=F+\dot m u$。
  \item 雨滴：$r=r_0+\lambda t$，$v=g/(4\lambda)[(r_0+\lambda t)-r_0^4/(r_0+\lambda t)^3]$。
  \item 碰撞链：$H=4h$（$m_2\ll m_1$）。
\end{enumerate}

\section{讨论}
\label{sec:discuss}

\paragraph{恢复系数的影响。}
$e=1$（完全弹性）时 $\beta_{\max}\to 90^\circ$（$\alpha\to0$）；
$e=0$（完全非弹性）时两球法向速度相同，但切向速度保留，$\beta$ 一般不为 $0$（只有正碰 $\alpha=0$ 才为 $0$）。

\paragraph{雨滴的终端行为。}
$v\sim gt/4$ 线性增长，无终端速度（因质量持续增加，重力加速度恒定）。

\paragraph{恢复系数的物理意义。}
$e$ 表征碰撞的``弹性程度''：$e=1$ 完全弹性（动能守恒），
$e=0$ 完全非弹性（动能最大损失）。实际碰撞 $0<e<1$。

\paragraph{变质量系统的应用。}
火箭推进（$dm/dt<0$）、雨滴增长、砂漏等都是变质量系统。
牛顿第二定律的推广形式 $\mathrm d(mv)/\mathrm dt=F+\dot m u$ 是关键。

\paragraph{与作业题的对照。}
表~\ref{tab:numeric} 覆盖了作业-5 Q1--Q3 的核心结论。

% ============================================================
\begin{thebibliography}{9}

\bibitem{goldstein}
H.~Goldstein, C.~P.~Poole, J.~L.~Safko,
\emph{Classical Mechanics}, 3rd ed., Pearson, 2002.
（碰撞与变质量系统。）

\bibitem{landau}
L.~D.~Landau, E.~M.~Lifshitz,
\emph{Mechanics}, 3rd ed., Butterworth--Heinemann, 1976.
（变质量系统。）

\bibitem{zhou}
周衍柏，《理论力学教程》（第三版），高等教育出版社，2009.
（碰撞与变质量。）

\end{thebibliography}

\end{document}
